\documentclass[11pt,a4paper]{article}

\usepackage[T1]{fontenc}
\usepackage[utf8]{inputenc}
\usepackage{lmodern}
\usepackage{microtype}
\usepackage[margin=24mm]{geometry}
\usepackage{amsmath,amssymb,mathtools}
\usepackage{booktabs,tabularx}
\usepackage{enumitem}
\usepackage{xcolor}
\usepackage{listings}
\usepackage{hyperref}
\usepackage{fancyhdr}

\definecolor{nismoblue}{HTML}{164E63}
\definecolor{nismogreen}{HTML}{166534}
\definecolor{nismored}{HTML}{991B1B}
\definecolor{nismogray}{HTML}{F3F4F6}
\definecolor{codegray}{HTML}{4B5563}

\hypersetup{
  colorlinks=true,
  linkcolor=nismoblue,
  urlcolor=nismoblue,
  citecolor=nismoblue,
  pdftitle={The Current Dynesty-Style rwalk Implementation in NISMO},
  pdfauthor={NISMO contributors}
}

\lstdefinestyle{nismopython}{
  language=Python,
  basicstyle=\ttfamily\small,
  keywordstyle=\color{nismoblue}\bfseries,
  commentstyle=\color{nismogreen},
  stringstyle=\color{nismored},
  numbers=left,
  numberstyle=\scriptsize\color{codegray},
  numbersep=8pt,
  backgroundcolor=\color{nismogray},
  frame=single,
  rulecolor=\color{gray!35},
  breaklines=true,
  showstringspaces=false,
  columns=fullflexible,
  keepspaces=true
}
\lstset{style=nismopython}

\newcommand{\NISMO}{\textsc{NISMO}}
\newcommand{\code}[1]{\texttt{\detokenize{#1}}}
\newcommand{\R}{\mathbb{R}}

\setlist{nosep,leftmargin=*}
\setlength{\parindent}{0pt}
\setlength{\parskip}{0.55em}
\renewcommand{\arraystretch}{1.18}

\pagestyle{fancy}
\fancyhf{}
\lhead{\NISMO{} standard \code{rwalk}}
\rhead{\thepage}
\setlength{\headheight}{14pt}

\title{\textbf{The Current Dynesty-Style \code{rwalk}\\
Implementation in \NISMO}}
\author{Implementation specification for \code{nismo} version
\code{0.1.0.dev3}}
\date{\today}

\begin{document}

\maketitle

\begin{abstract}
This document specifies the standard \code{proposal_scheme="rwalk"} kernel
currently implemented by \NISMO.  It follows the actual Python execution path,
including configuration resolution, construction and caching of the bounding
ellipsoid, uniform-ball proposals, the fixed-\(q_0\) Metropolis correction,
scale adaptation, tie handling, resource limits, and recorded diagnostics.
The controller and ellipsoid mechanics are adapted from Dynesty 3.1.0, while
the target density and acceptance rule are specific to \NISMO.
\end{abstract}

\tableofcontents
\newpage

\section{Scope and implementation locations}

This document concerns only the standard \code{"rwalk"} scheme.  The separate
\code{"en-rwalk"} differential-evolution ensemble kernel is not changed by
this implementation and retains its own covariance shrinkage and jitter
settings.

The implementation is divided among three modules:

\begin{itemize}
  \item \code{src/nismo/config.py}: public \code{RWalkSettings};
  \item \code{src/nismo/mcmc.py}: controller, ellipsoid, proposal, Metropolis
        transition, tuning, and method citation;
  \item \code{src/nismo/sampler.py}: run-level controller lifetime, dispatch,
        public citation property, live-set replacement, and histories.
\end{itemize}

No Dynesty package is imported at runtime.  The adapted source is distributed
under Dynesty's MIT terms in the repository's \code{NOTICE} file.

\section{The constrained density targeted by the walk}

Let \(L(\theta)\) be the likelihood, \(\pi(\theta)\) the normalized prior,
and \(q_0(\theta)\) the fixed normalized importance density supplied as
\code{importance_morph}.  \NISMO{} defines

\begin{equation}
  \log \Psi_0(\theta)
  = \log L(\theta)+\log\pi(\theta)-\log q_0(\theta).
\end{equation}

At one nested iteration, the discarded live point defines the threshold
\(\lambda\).  The replacement kernel must preserve

\begin{equation}
  p_\lambda(\theta)
  \propto q_0(\theta)
  \mathbf{1}\!\left[\log\Psi_0(\theta)>\lambda\right].
  \label{eq:target}
\end{equation}

This differs from Dynesty's usual unit-cube construction.  Dynesty obtains the
prior density through a prior transform, whereas \NISMO{} walks directly in
physical parameter space.  Consequently, merely passing the constraint is
not enough: \NISMO{} must also apply a Metropolis ratio involving \(q_0\).

\section{Public configuration and exact resolution}

The user passes an immutable \code{RWalkSettings} object:

\begin{lstlisting}
RWalkSettings(
    walks=None,
    facc=0.5,
    ncdim=None,
)
\end{lstlisting}

For model dimension \(D\), the run-level \code{RWalkSampler} resolves the
values as follows.

\begin{center}
\begin{tabularx}{\textwidth}{@{}l l X@{}}
\toprule
Field & Resolved value & Validation and meaning \\
\midrule
\code{walks}
& \(D+20\) when omitted
& An explicit value must be a positive integer and is raised to a minimum of
  two.  One successful replacement attempts exactly this many transitions. \\
\code{facc}
& \(\min(1,\max(1/W,f_{\rm input}))\)
& The input must be finite.  Here \(W\) is the resolved number of walks. \\
\code{ncdim}
& \(D\) when omitted
& An explicit value must equal \(D\).  Partial-dimensional refreshes cannot
  be implemented correctly for an arbitrary correlated importance Morph. \\
initial scale
& \(s_0=1\)
& Scale is controller state and is not a public setting. \\
\bottomrule
\end{tabularx}
\end{center}

The settings \code{n_steps}, \code{scale},
\code{covariance_shrinkage}, and \code{covariance_jitter} are not accepted by
standard \code{rwalk}.  Supplying them causes normal Python constructor errors.

Each call to \code{NISMOSampler.run()} constructs a fresh controller, so every
run begins with scale one, zero acceptance history, and no cached ellipsoid.

\section{Starting state and constraint ordering}

The current worst live point supplies \(\lambda\) but is never a chain state.
The implementation first identifies all surviving live indices that lie
strictly above the active constraint, removes the worst index explicitly, and
selects one eligible survivor uniformly.  If none exists, the replacement
fails with \code{insufficient_eligible_survivors}.

Under \code{tie_policy="strict"}, a proposal passes when

\begin{equation}
  \log\Psi_0(\theta')>\lambda.
\end{equation}

Under \code{tie_policy="randomized_plateau"}, every live state also contains
an auxiliary \(t\in[0,1)\).  The discarded point defines
\((\lambda,t_\lambda)\), every proposal draws a fresh \(t'\), and passage is
lexicographic:

\begin{equation}
  \log\Psi_0(\theta')>\lambda
  \quad\text{or}\quad
  \left[\log\Psi_0(\theta')=\lambda\ \text{and}\ t'>t_\lambda\right].
\end{equation}

A rejected transition retains both the current parameter vector and its
current auxiliary tie value.

\section{Dynesty-style single bounding ellipsoid}

\subsection{Input live set}

When no cached axes are available, or when the cache interval has elapsed,
the implementation constructs a single ellipsoid from the complete current
live matrix

\begin{equation}
  \Theta=(\theta_1,\ldots,\theta_N)^\mathsf{T}\in\R^{N\times D}.
\end{equation}

This includes the point defining the threshold.  Inclusion affects only the
proposal geometry; that point remains forbidden as a chain start.  At least
two finite live points and one dimension are required.

Let

\begin{equation}
  c=\frac{1}{N}\sum_{i=1}^{N}\theta_i,
  \qquad
  C=\operatorname{cov}(\Theta),
  \qquad
  \delta_i=\theta_i-c,
\end{equation}

where NumPy's sample covariance is used.  In one dimension the scalar result
is promoted to a \(1\times1\) matrix.

\subsection{Condition-number repair}

The covariance repair attempts at most 100 eigendecompositions.  Let
\(\ell_{\max}\) and \(\ell_{\min}\) be the largest and smallest eigenvalues.
The target maximum condition number is \(10^{12}\).

If the eigenvalues are finite and \(\ell_{\max}>0\), but

\begin{equation}
  \ell_{\min}<\frac{\ell_{\max}}{10^{12}},
\end{equation}

the eigenvalues are floored at

\begin{equation}
  \ell_j^{\rm fix}
  =\max\!\left(\ell_j,\frac{10\ell_{\max}}{10^{12}}\right),
\end{equation}

and the covariance is reconstructed.  If the spectrum is nonfinite, has no
positive eigenvalue, or eigendecomposition fails, the code blends toward the
identity:

\begin{equation}
  C\leftarrow(1-a_k)C+a_k I,
\end{equation}

where \(a_k\) grows geometrically from \(10^{-10}\) to one over the 100
attempts.  If no safe result is found, the code emits a runtime warning and
uses a unit sphere, \(C=P=A=I\).

For a successful spectrum \(C=V\operatorname{diag}(\ell)V^\mathsf{T}\), the
precision and proposal axes are

\begin{equation}
  P=V\operatorname{diag}(\ell^{-1})V^\mathsf{T},
  \qquad
  A=V\operatorname{diag}(\sqrt{\ell}).
\end{equation}

\subsection{Expansion to contain every live point}

The maximum squared ellipsoidal distance is

\begin{equation}
  d_{\max}=\max_i\delta_i^\mathsf{T}P\delta_i.
\end{equation}

If \(d_{\max}>0.999\), the implementation sets

\begin{equation}
  m=\frac{d_{\max}}{0.999},\qquad
  C\leftarrow mC,\quad
  P\leftarrow P/m,\quad
  A\leftarrow\sqrt{m}\,A.
\end{equation}

The repair-and-expansion process receives at most two passes.  Failure to
obtain an ellipsoid with every live point strictly inside unit ellipsoidal
distance raises \code{NumericalInvariantError}.

\subsection{Cache lifetime}

The axes are cached across replacements.  With \(W\) walks and \(N\) live
points, the refresh threshold is

\begin{equation}
  WN\quad\text{completed random-walk calls}.
\end{equation}

Each successful replacement adds exactly \(W\) to the cache counter.  Thus,
in the serial implementation, the bound is normally rebuilt before every
\(N\)-th subsequent replacement.  Interrupted or preflight-rejected chains do
not advance this controller counter because the run stops without completing
the replacement.

\section{One ellipsoidal-ball proposal}

For each transition, draw

\begin{equation}
  z\sim\mathcal{N}(0,I_D),qquad
  u\sim\operatorname{Uniform}(0,1),qquad
  r=u^{1/D}\frac{z}{\lVert z\rVert_2}.
\end{equation}

Then \(r\) is uniform in the unit \(D\)-ball and the proposal is

\begin{equation}
  \theta'=\theta+sAr,
  \label{eq:proposal}
\end{equation}

where \(A\) is the cached ellipsoid axis matrix and \(s\) is the scale frozen
at the start of the complete walk.  Equation~\eqref{eq:proposal} is symmetric
about the current state.  \NISMO{} does not clip, reflect, wrap, or redraw a
physical-space proposal.  Prior support is handled by ordinary evaluation and
rejection.

\section{Evaluation and Metropolis decision}

Every proposed parameter vector is evaluated exactly once by
\code{BatchEvaluator}, producing cached values

\begin{equation}
  \left(\theta',\log L',\log\pi',\log q_0',\log\Psi_0'\right).
\end{equation}

If the proposal fails the active pseudo-likelihood/tie constraint, it is
rejected without an MH random draw.  If it passes, the symmetric-proposal log
acceptance probability is

\begin{equation}
  \log\alpha
  =\min\left(0,\log q_0(\theta')-\log q_0(\theta)\right).
  \label{eq:mh}
\end{equation}

The implementation accepts when

\begin{equation}
  \log U < \log\alpha,
  \qquad U\sim\operatorname{Uniform}(0,1).
\end{equation}

An accepted transition replaces every cached field, including the tie value.
A rejected transition changes none of them.  A proposed \(\log q_0=-\infty\)
has zero acceptance probability; other nonfinite proposal densities raise a
numerical invariant error.  The eligible starting state must have finite
\(\log q_0\).

\section{Exact complete-walk algorithm}

For one nested replacement, the implemented sequence is:

\begin{enumerate}
  \item Validate live-array shapes and preflight the resource limits for all
        \(W\) proposals.
  \item Identify eligible survivors.  Fail if there are none.
  \item Obtain the cached or newly rebuilt ellipsoid axes from the complete
        live set.
  \item Select one eligible survivor uniformly and copy its complete cached
        state.
  \item Freeze the controller's current scale for this walk.
  \item Repeat exactly \(W\) times:
    \begin{enumerate}
      \item check the wall-time deadline;
      \item draw a uniform-ball displacement and form \(\theta'\);
      \item evaluate \(\theta'\) once and draw its tie value;
      \item test the active constraint;
      \item if valid, apply Equation~\eqref{eq:mh};
      \item on acceptance, replace the complete current state.
    \end{enumerate}
  \item Return the actual final state, even if no transition was accepted.
  \item Record \(W\) completed calls and tune the scale for the next nested
        replacement.
\end{enumerate}

The following compact pseudocode emphasizes the ordering:

\begin{lstlisting}
axes = sampler.axes_for(live_theta)
current = copy(uniform_choice(eligible_survivors))
scale_used = sampler.scale

for _ in range(sampler.walks):
    r = uniform_unit_ball(ndim)
    proposed = evaluate(current.theta + scale_used * axes @ r)
    proposed.tie = rng.random()

    if not passes_constraint(proposed):
        continue
    n_valid += 1

    log_alpha = min(0.0, proposed.log_q0 - current.log_q0)
    if log(rng.random()) < log_alpha:
        current = proposed
        n_accepted += 1

sampler.record_completed_walk(
    accept=n_accepted,
    scale=scale_used,
)
return current
\end{lstlisting}

\section{Scale adaptation}

Let \(A_k\) and \(R_k\) denote accepted and rejected transitions accumulated
in the controller.  In the current serial \NISMO{} path, tuning updates after
every completed replacement, so normally

\begin{equation}
  A_k+R_k=W,
  \qquad
  f_k=\frac{A_k}{A_k+R_k}.
\end{equation}

Here ``accepted'' means final MH acceptance.  Constraint failures and valid
but MH-rejected proposals both contribute to \(R_k\).  With target fraction
\(f_0=\code{facc}\), the Dynesty recursion is

\begin{equation}
  s_{k+1}
  =s_k\exp\left(\frac{f_k-f_0}{D f_0}\right).
  \label{eq:tune}
\end{equation}

Acceptance above target expands the next walk; acceptance below target shrinks
it.  After Equation~\eqref{eq:tune}, both accumulated history counters reset
to zero.  The optional \code{update=False} controller path only accumulates
history and is retained from the Dynesty-style interface, but the current
serial \NISMO{} dispatch updates immediately after a completed walk.

\section{Resource limits and failure semantics}

Before constructing a bound or starting a chain, \NISMO{} rejects the complete
replacement without any new likelihood calls when:

\begin{itemize}
  \item \(W>\code{max_proposals_per_replacement}\);
  \item the run-wide likelihood budget cannot accommodate all \(W\) calls;
  \item the wall-time deadline has already elapsed.
\end{itemize}

The wall-time deadline is also checked before every scalar transition.  If it
expires inside a walk, the replacement returns no draw, the worst live point
is left untouched, and the shortened chain is never used.  A successful
replacement always has

\begin{equation}
  n_{\rm proposed}=n_{\rm likelihood\ calls}=n_{\rm completed}=W.
\end{equation}

This equality applies to calls made by the replacement itself; initial live
point evaluation is additional run-wide work.

\section{Recorded diagnostics}

For each completed nested replacement, \code{RunHistory} records:

\begin{center}
\begin{tabularx}{\textwidth}{@{}l X@{}}
\toprule
Field & Standard \code{rwalk} meaning \\
\midrule
\code{constraint_pass_fraction}
& Number passing the active pseudo-likelihood/tie constraint divided by \(W\). \\
\code{mh_acceptance_fraction}
& Number accepted by the complete constraint-plus-MH transition divided by \(W\). \\
\code{mcmc_accepted}
& Count of accepted transitions in the completed walk. \\
\code{mcmc_moved}
& One if at least one transition was accepted, otherwise zero. \\
\code{mcmc_completed}
& Exactly \(W\). \\
\code{acceptance_fraction}
& Cumulative nested replacements divided by all proposals; this retains the
  sampler-wide efficiency meaning and is not the MH rate. \\
\bottomrule
\end{tabularx}
\end{center}

The adapted scale and ellipsoid axes are controller state and are not added as
public \code{RunHistory} arrays.

\section{Why the kernel preserves the desired density}

During one complete walk, \(A\) and \(s\) are frozen.  The uniform-ball
increment is centered and symmetric, hence

\begin{equation}
  Q(\theta'\mid\theta)=Q(\theta\mid\theta').
\end{equation}

Inside the permitted constraint, the target-density ratio from
Equation~\eqref{eq:target} is exactly

\begin{equation}
  \frac{p_\lambda(\theta')}{p_\lambda(\theta)}
  =\frac{q_0(\theta')}{q_0(\theta)},
\end{equation}

which gives Equation~\eqref{eq:mh}.  Proposals outside the constraint have
target density zero and are rejected.  Bound and scale adaptation occurs only
between complete walks, so each individual transition kernel retains the
required symmetry.

Invariance does not imply independence.  The returned point remains
correlated with its selected live survivor, especially when \(W\) is too
small or the acceptance rate is poor.  Users should compare repeated runs,
monitor the separate constraint and MH fractions, and test sensitivity to
\code{walks} and \code{n_live}.  Conditioning the output on movement, forcing
an acceptance, or choosing the highest-likelihood visited state would change
the target and is not implemented.

\section{Usage}

Explicit control:

\begin{lstlisting}
from nismo import NISMOSampler, RWalkSettings

sampler = NISMOSampler(
    model=model,
    importance_morph=importance_morph,
    proposal_scheme="rwalk",
    rwalk_settings=RWalkSettings(
        walks=50,
        facc=0.5,
    ),
    n_live=200,
    rng=42,
)

result = sampler.run(
    dlogz=0.1,
    max_proposals_per_replacement=100,
)
\end{lstlisting}

Dynesty-style top-level defaults:

\begin{lstlisting}
sampler = NISMOSampler(
    model=model,
    importance_morph=importance_morph,
    proposal_scheme="rwalk",
    rwalk_settings=RWalkSettings(),
    n_live=200,
    rng=42,
)
\end{lstlisting}

For this second example, \(W=D+20\), \(f_0=0.5\), and \(s_0=1\).  The hard
proposal limit supplied to \code{run()} must be at least \(W\).

Useful inspection:

\begin{lstlisting}
print(sampler.citations)
print(result.history.constraint_pass_fraction)
print(result.history.mh_acceptance_fraction)
print(result.history.mcmc_moved)
\end{lstlisting}

For standard \code{rwalk}, \code{sampler.citations} returns only the Skilling
(2006) method entry, as specified by the current public behavior.

\section{Attribution}

The stateful controller, uniform-unit-ball draw, covariance repair, and
single-ellipsoid construction are adapted from the Dynesty 3.1.0 source:

\begin{center}
\url{https://github.com/joshspeagle/dynesty/tree/v3.1.0}
\end{center}

Dynesty's copyright and MIT permission notice are reproduced in the root
\code{NOTICE} file and included in source and wheel distributions.  The
runtime method citation remains the reference below.

\begin{thebibliography}{9}

\bibitem{skilling2006}
J. Skilling,
``Nested sampling for general Bayesian computation,''
\emph{Bayesian Analysis}, vol.~1, no.~4, pp.~833--859, 2006.
\url{https://projecteuclid.org/euclid.ba/1340370944}

\end{thebibliography}

\end{document}
